% Generated draft: requires mathematical and source review.
\documentclass[11pt]{article}
\usepackage[margin=25mm]{geometry}
\usepackage{fontspec}
\setmainfont{Latin Modern Roman}
\usepackage{amsmath,amssymb,mathtools,unicode-math}
\setmathfont{Latin Modern Math}
\usepackage{enumitem,needspace,xurl,hyperref}
\setlength{\parindent}{0pt}
\setlength{\parskip}{5pt}
\setlength{\emergencystretch}{3em}
\newcommand{\N}{\mathbb{N}}
\newcommand{\Z}{\mathbb{Z}}
\newcommand{\Q}{\mathbb{Q}}
\newcommand{\R}{\mathbb{R}}
\newcommand{\C}{\mathbb{C}}
\newcommand{\F}{\mathbb{F}}
\newcommand{\A}{\mathbb{A}}
\newcommand{\PP}{\mathbb{P}}
\newcommand{\E}{\mathbb{E}}
\newcommand{\eps}{\varepsilon}
\newcommand{\dd}{\,\mathrm{d}}
\begin{document}
\section*{UnsolvedMath ID: 11000283}
\textbf{Source problem code:} catalogue entry
\subsection{Definitions and mathematical statement}
Outer space in the usual rank-\(n\geq2\) model consists of marked finite core metric graphs of rank \(n\), with positive edge lengths and total volume \(1\), and carries the \(\operatorname{Out}(F_n)\)-action. A standard orbit-walk setup chooses a probability measure \(\mu\) on \(\operatorname{Out}(F_n)\), independent increments \(\phi_i\) with that law, products \(\Phi_k=\phi_1\cdots\phi_k\), and orbit points \(\Phi_kx_0\) in Outer space.

The question asks for a random-walk theory analogous to Kaimanovich--Masur's Teichm\"uller-space theory. The excerpt does not specify the probability law, moment/entropy or non-elementarity assumptions, metric, compactification or which analogous conclusions are sought. Boundary convergence, hitting measures or a Poisson-boundary identification would each need precise additional hypotheses; none is asserted for all measures merely from this analogy.

The following sentence suggests bounded cohomology, modeled on Bestvina--Fujiwara's mapping-class subgroup work [6], as an approach to ``Question 17.'' That numbered question is absent and is not identified with the present one, so the remark is retained without inventing a bounded-cohomology theorem or a dependency. The supplied summary reports developed random-outer-automorphism theory and mentions Rivin and others, but supplies no exact comparison theorem or complete reference.
\subsection{Short English statement}
Develop an Outer-space analogue of the intended Kaimanovich--Masur random-walk theory, with probability, boundary and conclusion hypotheses specified; the separate Question 17 reference remains unresolved.
\subsection{Sources}
\begin{itemize}
\item \emph{Problems on Mapping Class Groups and Related Topics}, edited by Farb, Question 3.3, supplied PDF page 333. \url{https://www.math.uchicago.edu/~farb/papers/mcgbook.pdf}. Recorded access 2026-07-29; extraction pdf-font-structure.
\item Kaimanovich and Masur, named Teichm\"uller-space analogy; Bestvina and Fujiwara [6], bounded-cohomology/subgroup template. Full bibliography and Question 17 are absent. The summary additionally names Rivin and others without complete citations.
\item Local UnsolvedMath registry, \texttt{data/unsolvedmath/problems.json}, permanent ID \(11000283\), exact source code \texttt{catalogue entry}; supplied OPEN-TRIAGE summary reports substantial related progress.
\end{itemize}
\end{document}
